\section{Introduction}
Neural receivers (NRXs) replace channel estimation, equalization, and demapping with a single neural network trained on data~\cite{honkala2021deeprx}. Cammerer et al.~\cite{cammerer2023neural} brought this concept to multi-user multiple-input multiple-output (MU-MIMO) transmission over the 5G NR physical uplink shared channel (PUSCH), using convolutional and graph neural network (CGNN) layers that refine a state per user over several iterations. More iterations yield more reliable log-likelihood ratios (LLRs), but each iteration also adds latency. In the real-time implementation of this NRX, the 1\,ms budget admits only two of its eight iterations~\cite{wiesmayr2025realtime}. On our evaluation sequences, the cut to two iterations clearly raises the signal-to-noise ratio (SNR) required for a transport block error rate (TBLER) of 10\%, even when the NRX is retrained for two iterations.

Like other NRXs, this receiver processes each slot independently. Yet the channel of a user evolves continuously, so the pilots of earlier slots carry information about the current channel, which the receiver discards with each slot. Classical receivers exploit the temporal correlation of the channel by filtering pilots across time or by bundling the demodulation reference signals (DMRS) of several slots~\cite{hoeher1997wiener,3gpp38214}.

In this paper, we investigate whether an NRX can also leverage the pilots of earlier slots to substitute for iterations. We propose a stateful NRX that carries channel estimates across slots. It runs two CGNN iterations and keeps, for each user, a pilot memory with the least-squares (LS) channel estimates at the DMRS positions of up to seven previous transport blocks (TBs). A channel filter combines these entries with the current LS estimate before the first iteration. The filter coefficients are estimated online for each user by minimizing the prediction error of the historical channel estimates in the pilot memory, without prior channel statistics or ground-truth channel labels. The filter also selects how many historical estimates to use by how well they predict the current pilots, which mitigates outdated channel information after channel changes or scheduling gaps. We train the NRX on sequences of TBs rather than independent slots, starting from the released eight-iteration NRX, and add no trained weights. With two iterations, it recovers most of the accuracy that the cut from eight iterations costs.

In summary, our contributions are:
\begin{itemize}
\item We develop a stateful NRX with a pilot memory and a channel filter fitted on it without training.
\item We evaluate the proposed NRX in simulations of $2\times 4$ MU-MIMO with QPSK, 16-QAM, and 64-QAM over sequences of 64 slots, with and without power control: it clearly outperforms a stateless two-iteration NRX retrained on the same data and the released NRXs with up to three iterations, and performs close to the eight-iteration NRX.
\item It adds no CGNN weights, and our latency measurements with TensorRT show that it remains real-time capable, unlike the released three-iteration NRX.
\end{itemize}
